
\documentclass{article}
\usepackage{colortbl}
\usepackage{makecell}
\usepackage{multirow}
\usepackage{supertabular}

\begin{document}

\newcounter{utterance}

\centering \large Interaction Transcript for game `ifeval', experiment `ifeval', episode 62 with Qwen3.5{-}4B.
\vspace{24pt}

{ \footnotesize  \setcounter{utterance}{1}
\setlength{\tabcolsep}{0pt}
\begin{supertabular}{c@{$\;$}|p{.15\linewidth}@{}p{.15\linewidth}p{.15\linewidth}p{.15\linewidth}p{.15\linewidth}p{.15\linewidth}}
    \# & \multicolumn{2}{c}{Player} && \multicolumn{2}{c}{Game Master} \\
    \hline

    \theutterance \stepcounter{utterance}  
    & & & \multicolumn{4}{p{0.6\linewidth}}{
        \cellcolor[rgb]{0.9,0.9,0.9}{
            \makecell[{{p{\linewidth}}}]{
                \texttt{\tiny{[P1$\langle$GM]}}
                \texttt{Which of the following is not a fish: salmon or avocado? Answer this easy question first, then expand into an interesting article about salmon and avocado. Use some words in all caps in your response, but those all{-}caps words should appear at most 10 times.} \\
            }
        }
    }
     \\ \\

    \theutterance \stepcounter{utterance}  
    & \multicolumn{4}{p{0.6\linewidth}}{
        \cellcolor[rgb]{0.561,0.737,0.561}{
            \makecell[{{p{\linewidth}}}]{
                \texttt{\tiny{[P1$\rangle$GM]}}
                \texttt{The answer is simple: **AVOCADO** is not a fish; salmon is a fish.} \\
\\ 
\texttt{\#\#\# The Perfect Pair: Salmon and Avocado} \\
\\ 
\texttt{In the culinary world, the pairing of **SALMON** and **AVOCADO** is often hailed as a classic combination, yet their biological origins are worlds apart. Salmon belongs to the family Salmonidae, swimming through cold, clear waters of the Pacific and Atlantic oceans. They are renowned for their rich, oily flesh, which is packed with Omega{-}3 fatty acids essential for heart health. In contrast, the avocado is a fruit native to Central and South America, specifically the genus *Persea*. Despite its green skin and creamy texture, it has absolutely no relation to aquatic life.} \\
\\ 
\texttt{When you place a slice of ripe avocado atop a seared salmon fillet, you are creating a textural and flavor symphony. The **SMOOTH** creaminess of the avocado contrasts beautifully with the **FIRM** flaky texture of the fish. The mild, buttery taste of the avocado complements the distinct, slightly sweet flavor of the salmon without overpowering it. This dish is not just delicious; it is a nutritional powerhouse. The combination provides a perfect balance of healthy fats, high{-}quality protein, and fiber.} \\
\\ 
\texttt{While one swims and the other grows in the soil, their union on your plate proves that nature's best ingredients often come from unexpected places. Whether enjoyed in a fresh sashimi style or grilled with a zesty lemon dressing, this duo remains a staple in modern gastronomy. It is a reminder that while we explore the depths of the ocean for **SALMON**, we look to the earth for the **AVOCADO**, yet both bring us closer to a healthier, more flavorful life.} \\
            }
        }
    }
    & & \\ \\

    \theutterance \stepcounter{utterance}  
    & & & \multicolumn{2}{p{0.3\linewidth}}{
        \cellcolor[rgb]{0.9,0.9,0.9}{
            \makecell[{{p{\linewidth}}}]{
                \texttt{\tiny{[GM$|$GM]}}
                {'change_case:capital_word_frequency': {'capital_relation': 'at least', 'capital_frequency': 1}}
            }
        }
    }
    & & \\ \\

    \theutterance \stepcounter{utterance}  
    & & & \multicolumn{2}{p{0.3\linewidth}}{
        \cellcolor[rgb]{0.9,0.9,0.9}{
            \makecell[{{p{\linewidth}}}]{
                \texttt{\tiny{[GM$|$GM]}}
                \texttt{game\_result = WIN} \\
            }
        }
    }
    & & \\ \\

\end{supertabular}
}

\end{document}
